% content/indonesian/glossary.tex
% -----------------------------------------------------------------------------
% GLOSARIUM
% -----------------------------------------------------------------------------
% Muat singkatan beserta kepanjangannya dan istilah khusus bidang yang dipakai
% dalam naskah. Istilah yang sudah umum dikenal pembaca program studi tidak
% perlu dicantumkan. Entri diurutkan menurut abjad.
%
% Judul halaman ditulis oleh content/frontmatter-lists.tex, jadi berkas ini
% cukup berisi tabelnya saja. Pada naskah proposal dan draf, glosarium ini
% sekaligus menggantikan DAFTAR SINGKATAN.

\TemplateTodo{Ganti tiga entri di bawah ini dengan istilah dan singkatan yang benar-benar dipakai pada naskah Anda.}

\begin{longtable}{@{}p{0.26\linewidth}p{0.68\linewidth}@{}}
API & \textit{Application Programming Interface}, antarmuka yang memungkinkan satu perangkat lunak memanggil layanan perangkat lunak lain. \\
CNN & \textit{Convolutional Neural Network}, jaringan saraf tiruan yang memakai operasi konvolusi untuk mengekstraksi ciri spasial dari citra. \\
Prapemrosesan & Rangkaian langkah penyiapan data mentah sebelum data dipakai untuk pelatihan model, misalnya penyeragaman ukuran dan normalisasi. \\
\end{longtable}
